\subsubsection{Service}

Kubernetes baut ein eigenes internes Netzwerk auf, über das die Pods miteinander kommunizieren können und erreichbar sind.
Dabei erhält jeder Pod eine eigene IP-Adresse, allerdings ist ein Pod kurzlebig. Wenn ein Pod ausfällt und ersetzt wird, hat der neue Pod nicht die gleiche IP-Adresse wie der Pod davor.

Damit die Pods mit einer festen IP-Adresse erreicht werden können, gibt es in Kubernetes Services, die alle Replicas eines Pods kennen.
Zusätzlich dient ein Service als Loadbalancer, der den Datenverkehr auf die Pods verteilt.
Die Kommunikation findet daher nicht über die Pods selbst, sondern über deren Service statt.